\chapter{Validación de la aplicación}\label{chapter:result}


% \section{Propósito de la validación}

% La validación del sistema persigue demostrar cuatro afirmaciones concretas:

% \begin{enumerate}
%     \item \textbf{Los flujos principales operan correctamente.} El registro de reportes, la asignación a revisores, la gestión de revisiones y las notificaciones automáticas funcionan conforme a los requerimientos funcionales definidos en el Capítulo \ref{chapter:proposal}.
%     \item \textbf{El sistema responde en menos de 5 segundos.} El envío de un reporte completo —desde la recepción de la solicitud hasta la respuesta de confirmación— no supera el límite establecido en el RNF-05, excluyendo la latencia de red externa.
%     \item \textbf{Los mecanismos de seguridad operan según lo diseñado.} Los endpoints protegidos rechazan solicitudes sin token, con token expirado o con rol insuficiente; el rate limiting bloquea el exceso de solicitudes; y el CAPTCHA impide el envío automatizado de formularios.
%     \item \textbf{Los usuarios finales consideran la interfaz útil y fácil de usar.} La evaluación mediante el cuestionario SUS (System Usability Scale) debe arrojar una puntuación superior a 68, umbral que define la usabilidad aceptable según el estándar del cuestionario \cite{brooke1996sus}.
% \end{enumerate}

% \section{Entorno y metodología}

% \subsection{Hardware}

% Las pruebas se ejecutaron sobre una máquina con las siguientes características:

% \begin{itemize}
%     \item \textbf{Procesador:} Intel Core i5-8250 CPU @ 1.80GHz
%     \item \textbf{Memoria RAM:} 16 GB DDR4
%     \item \textbf{Almacenamiento:} SSD KINGSTON de 512 GB
% \end{itemize}

% \subsection{Software}

% El entorno de pruebas replicó la configuración de desarrollo local definida en el Capítulo \ref{chapter:implementation}:

% \begin{itemize}
%     \item \textbf{Sistema operativo:} Windows 11 Pro

%     \item \textbf{Base de datos:} MySQL 8.0
%     \item \textbf{Backend:} .NET 8, ejecutado mediante Docker Compose
%     \item \textbf{Frontend:} Vite + React 18, ejecutado en modo desarrollo
%     \item \textbf{Navegador:} Google Chrome 124
% \end{itemize}

% \subsection{Red}

% Todas las pruebas se ejecutaron en entorno local, con la API y la base de datos corriendo en la misma máquina. Esta configuración excluye la latencia de red externa de las mediciones de rendimiento, conforme al RNF-05.

% \subsection{Datos de prueba}

% La base de datos se pobló con datos sintéticos que simulan el escenario real del Instituto Finlay:

% \begin{itemize}
%     \item \textbf{Provincias y municipios:} 16 provincias y 32 municipios, correspondientes a la división político-administrativa de Cuba.
%     \item \textbf{Vacunas:} 5 vacunas activas con sus respectivos fabricantes.
%     \item \textbf{Síntomas:} 15 síntomas activos.
%     \item \textbf{Usuarios:} 1 administrador, 2 jefes de vacunación municipales, 3 médicos evaluadores y 2 ciudadanos de prueba.
%     \item \textbf{Reportes:} 20 reportes de ESAVI con distintos estados, severidades y asignaciones.
% \end{itemize}

% \subsection{Metodología}

% La validación combinó cuatro enfoques complementarios:

% \begin{itemize}
%     \item \textbf{Pruebas funcionales manuales:} ejecución de los flujos principales del sistema por parte del desarrollador, verificando que cada paso produce el resultado esperado según los requerimientos funcionales.
%     \item \textbf{Pruebas unitarias automatizadas:} validación de reglas de negocio mediante xUnit sobre los validadores de la capa de aplicación, con dependencias simuladas.
%     \item \textbf{Pruebas de rendimiento:} medición de tiempos de respuesta con K6, ejecutando 100 solicitudes por endpoint.
%     \item \textbf{Pruebas de usabilidad:} evaluación con usuarios representativos mediante el método de pensamiento en voz alta y el cuestionario SUS.
% \end{itemize}

\section{Pruebas funcionales}

Las pruebas funcionales verificaron que los flujos principales del sistema operan conforme a los requerimientos funcionales definidos en el Capítulo \ref{chapter:proposal}. La Tabla \ref{tab:pruebas_funcionales} resume los casos de prueba ejecutados, seleccionados por su representatividad dentro del flujo de farmacovigilancia.

\begin{table}[H]
\centering
\caption{Resumen de casos de prueba funcionales}
\label{tab:pruebas_funcionales}
\begin{tabular}{|p{1.2cm}|p{3.5cm}|p{3.5cm}|p{3cm}|c|}
\hline
\textbf{ID} & \textbf{Caso de prueba} & \textbf{Datos de entrada} & \textbf{Resultado esperado} & \textbf{¿Pasó?} \\
\hline
CP01 & Registro de reporte público & Datos completos de sujeto, vacuna y síntoma & Código 201 y número de notificación & \checkmark \\
\hline
CP02 & Registro con campos vacíos & Formulario sin datos obligatorios & Código 400 con errores de validación & \checkmark \\
\hline
CP03 & Inicio de sesión con credenciales válidas & Correo y contraseña correctos & JWT, cookie de refresh y redirección al panel & \checkmark \\
\hline
CP04 & Inicio de sesión con credenciales inválidas & Correo y contraseña incorrectos & Código 400 con mensaje de error & \checkmark \\
\hline
CP05 & Asignación de reporte a revisor & Reporte pendiente y revisor disponible & Asignación creada y notificación enviada & \checkmark \\
\hline
CP06 & Consulta de reportes por municipio & Municipio con 5 reportes registrados & Lista con 5 reportes & \checkmark \\
\hline
CP07 & Finalización de revisión médica & Reporte en estado: \textit{En revisión} & Estado actualizado a \texttt{Revisado} & \checkmark \\
\hline
CP08 & Exportación de reporte a PDF & Reporte con datos completos & Archivo PDF generado correctamente & \checkmark \\
\hline
CP09 & Consultar estado de reporte válido & Número de notificación & Estado detallado & \checkmark \\
\hline
\end{tabular}
\end{table}


\section{Pruebas de rendimiento}

\section{Pruebas de usabilidad}

\section{Discusión de resultados}

\section{Resumen de hallazgos}